\chapter{Empalme torsional de una región atrapada y construcción de una rama cosmológica descendiente}
\label{chap:universo-descendiente}

\section{Empalme torsional de una región atrapada: estratos exacto, condicionado y físico de la rama cosmológica descendiente}

El \cref{ch:rebote-ec} establece la ley de memoria
\begin{equation}
 M_n=\left(\frac{a_n}{a_0}\right)^{-6}
 \label{eq:descendant-memory-law}
\end{equation}
y formula su realizaci\'on mediante una densidad cuadr\'atica de esp\'in en
Einstein--Cartan. El presente cap\'itulo estudia el paso adicional desde un
rebote torsional hacia una regi\'on causal expansiva situada tras el interior
atrapado de un agujero negro.

La exposici\'on distingue \(3\) estratos:
\begin{enumerate}
 \item \emph{estructura exacta}: orientaci\'on TRIT, involuci\'on de M\"obius,
       conservaci\'on isom\'etrica del canal y ley \(V^{-2}\);
 \item \emph{consecuencia condicionada}: rebote de Einstein--Cartan,
       empalme geom\'etrico y apertura causal bajo hip\'otesis declaradas;
 \item \emph{realizaci\'on geom\'etrica clasificada}: existencia explícita de
       la continuaci\'on descendiente m\'inima y obstrucci\'on del cobordismo
       ramificado cl\'asico con cambio topol\'ogico.
\end{enumerate}
Esta estratificaci\'on permite desarrollar la imagen completa y conservar el
estatuto de cada afirmaci\'on.

El antecedente com\'un es un estado enriquecido
\begin{equation}
 X^{\rm enr}_{\TPK}
 =(x,\text{fibra},\text{orientaci\'on},\text{carry},
   \text{frontera},\text{ruta},\text{memoria},\text{observador}),
 \label{eq:descendant-enriched-state}
\end{equation}
producido por \(\mathrm{APP}\to\TRIT\to\TPK\). Los lectores causal,
torsional, topol\'ogico y cu\'antico se aplican a este mismo estado y
conservan su genealog\'ia.

\section{Orientación causal de las \(3\) hojas}

La gram\'atica de TRIT se realiza mediante las \'algebras
\begin{equation}
 \mathbb A_\tau=\R[X]/(X^2+\tau),
 \qquad
 J_\tau^2=-\tau I,
 \qquad
 \tau\in\{+1,0,-1\}.
 \label{eq:descendant-TRIT-algebras}
\end{equation}
Sus exponenciales son
\begin{align}
 \ee^{tJ_{+1}}&=\cos t\,I+\sin t\,J_{+1},
 \label{eq:descendant-elliptic}\\
 \ee^{tJ_0}&=I+tJ_0,
 \label{eq:descendant-parabolic}\\
 \ee^{tJ_{-1}}&=\cosh t\,I+\sinh t\,J_{-1}.
 \label{eq:descendant-hyperbolic}
\end{align}
La hoja \(+1\) porta cierre, retorno y memoria; la hoja \(0\) porta umbral,
presente y primera variaci\'on; la hoja \(-1\) porta apertura, propagaci\'on y
futuro two-way.

Sea \(\vartheta\) un escalar de expansi\'on orientado. Definimos el lector
causal
\begin{equation}
 \boxed{
 \tau_{\rm cau}(\vartheta):=-\operatorname{sgn}(\vartheta)
 \in\{-1,0,+1\}.}
 \label{eq:descendant-causal-reader}
\end{equation}
Una expansi\'on positiva publica apertura \(-1\); una expansi\'on nula
publica umbral \(0\); una expansi\'on negativa publica cierre \(+1\).

En una familia de esferas, \(\vartheta=\theta_{(\ell)}\) puede tomarse como
la expansi\'on nula saliente. En una regi\'on homog\'enea o anis\'otropa,
\(\vartheta=\Theta\) puede tomarse como la expansi\'on volum\'etrica. Estos
dos observables ocupan dominios distintos y comparten la misma regla de
signo.

\begin{proposition}[Secuencia causal de doble umbral]
\label{prop:descendant-double-threshold}
Sup\'onganse las condiciones de signo
\begin{equation}
 \theta_{(\ell)}>0\ \text{en el exterior},
 \quad
 \theta_{(\ell)}=0\ \text{en el horizonte},
 \quad
 \theta_{(\ell)}<0\ \text{en el interior atrapado},
 \label{eq:descendant-null-signs}
\end{equation}
y
\begin{equation}
 \Theta<0\ \text{antes del rebote},
 \quad
 \Theta=0\ \text{en el rebote},
 \quad
 \Theta>0\ \text{despu\'es del rebote}.
 \label{eq:descendant-volume-signs}
\end{equation}
El lector \cref{eq:descendant-causal-reader} produce la secuencia
\begin{equation}
 \boxed{
 -1_{\rm exterior}
 \longrightarrow0_{\rm horizonte}
 \longrightarrow+1_{\rm interior}
 \longrightarrow0_{\rm rebote}
 \longrightarrow-1_{\rm descendiente}.}
 \label{eq:descendant-five-stage-sequence}
\end{equation}
\end{proposition}

\begin{proof}
La afirmaci\'on resulta de aplicar
\(\tau_{\rm cau}(\vartheta)=-\operatorname{sgn}(\vartheta)\) a cada signo de
\cref{eq:descendant-null-signs,eq:descendant-volume-signs}.
\end{proof}

Los dos estados \(0\) son realizaciones del mismo arquetipo de umbral. El
primero expresa marginalidad nula; el segundo expresa marginalidad
volum\'etrica. La identidad del signo conserva la gram\'atica, mientras el
observable declarado determina el contenido f\'isico de cada umbral.

\section{Interior anis\'otropo y memoria cuadr\'atica}

Una regi\'on interior esf\'ericamente sim\'etrica admite la reducci\'on de
Kantowski--Sachs
\begin{equation}
 \dd s^2=-N(\tau)^2\dd\tau^2
 +a_\parallel(\tau)^2\dd\chi^2
 +a_\perp(\tau)^2\dd\Omega_2^2.
 \label{eq:descendant-KS-metric}
\end{equation}
El volumen de una celda com\'ovil es
\begin{equation}
 V(\tau)=V_0a_\parallel a_\perp^2,
 \label{eq:descendant-KS-volume}
\end{equation}
y su expansi\'on es
\begin{equation}
 \Theta
 :=\frac1N\frac{\dot V}{V}
 =H_\parallel+2H_\perp,
 \qquad
 H_i:=\frac1N\frac{\dot a_i}{a_i}.
 \label{eq:descendant-KS-expansion}
\end{equation}

Una carga com\'ovil de esp\'in \(N_s\) produce una densidad
\begin{equation}
 n_s=\frac{N_s}{V},
 \label{eq:descendant-spin-number-density}
\end{equation}
y la ecuaci\'on de Cartan, lineal en la corriente, genera una correcci\'on
cuadr\'atica
\begin{equation}
 \rho_s=s_0\left(\frac V{V_0}\right)^{-2}.
 \label{eq:descendant-spin-density}
\end{equation}
La rama de Perron publica exactamente la misma covariancia:
\begin{equation}
 M_n=V_n^{-2}.
 \label{eq:descendant-Perron-volume}
\end{equation}
El transductor torsional
\begin{equation}
 \mathcal T_{\rm EC}:M_n\longmapsto
 \rho_{s,n}=s_0M_n
 \label{eq:descendant-EC-transducer}
\end{equation}
conserva positividad y el exponente de volumen. La amplitud \(s_0\) procede
de la corriente de Cartan y mantiene su obligaci\'on de selecci\'on
prospectiva.

Para la discusi\'on local adoptamos unidades geom\'etricas \(c=1\). Una
restricci\'on efectiva de Kantowski--Sachs tiene la forma
\begin{equation}
 \frac{\Theta^2}{3}
 =8\pi G(\rho_m-\rho_s)
 +\sigma^2-\frac12\,{}^{(3)}R+\Lambda,
 \label{eq:descendant-KS-constraint}
\end{equation}
donde \(\sigma^2\) es la cizalla y \({}^{(3)}R\) la curvatura de la hoja.
Un rebote volum\'etrico regular satisface
\begin{equation}
 \boxed{
 \Theta(\tau_b)=0,
 \qquad
 \frac1N\dot\Theta(\tau_b)>0,}
 \label{eq:descendant-bounce-conditions}
\end{equation}
junto con finitud positiva de \(a_\parallel,a_\perp\) y de los invariantes
de curvatura y torsi\'on.

\begin{criterion}[Dominancia torsional anis\'otropa]
\label{crit:descendant-torsion-dominance}
Sup\'ongase que, durante la contracci\'on,
\begin{equation}
 \rho_s=s_0V^{-2},
 \qquad
 \sigma^2=\Sigma_0^2V^{-2}+o(V^{-2}),
 \label{eq:descendant-torsion-shear-scaling}
\end{equation}
y que materia, curvatura y \(\Lambda\) crecen con potencia estrictamente
menor que \(V^{-2}\). La desigualdad
\begin{equation}
 \boxed{8\pi Gs_0>\Sigma_0^2}
 \label{eq:descendant-torsion-dominates-shear}
\end{equation}
es una condici\'on suficiente para que el t\'ermino torsional domine la
cizalla en el r\'egimen de volumen peque\~no. La existencia de una ra\'iz
transversal de \cref{eq:descendant-KS-constraint} y la desigualdad
\cref{eq:descendant-bounce-conditions} completan el criterio de rebote.
\end{criterion}

La dominancia asint\'otica aporta el mecanismo; la ra\'iz transversal aporta
el evento geom\'etrico. Esta distinci\'on conserva la dependencia respecto de
curvatura, anisotrop\'ia, ecuaci\'on de estado y datos de frontera.

\section{Condici\'on non\'adica para el nivel de rebote}

Escribamos \(x:=a/a_0\) para la escala is\'otropa adimensional. En los
retornos completos \(n=9q\), la rama de Perron satisface
\begin{equation}
 x_{9q}=\frac{a_{9q}}{a_0}=\varphi_{\HMT}^{3q}.
 \label{eq:descendant-nonadic-scale}
\end{equation}
En la reducci\'on is\'otropa con
\begin{equation}
 \rho_m=\rho_0x^{-3(1+w)},
 \qquad
 \rho_s=s_0x^{-6},
 \qquad
 w<1,
 \label{eq:descendant-isotropic-densities}
\end{equation}
la igualdad de densidades fija
\begin{equation}
 x_b^{3(1-w)}=\frac{s_0}{\rho_0}.
 \label{eq:descendant-isotropic-bounce}
\end{equation}
Si el rebote se exige sobre una vuelta completa del TPK, la eliminaci\'on de
\(x_b\) produce
\begin{equation}
 \boxed{
 q_b=
 \frac{\log(s_0/\rho_0)}
 {9(1-w)\log\varphi_{\HMT}},
 \qquad
 q_b\in\Z.}
 \label{eq:descendant-quantized-bounce}
\end{equation}
La integridad de \(q_b\) es una condici\'on adicional de sincronizaci\'on,
derivada de la hip\'otesis de coincidencia entre el rebote y el retorno
non\'adico. Un rebote continuo entre dos vueltas satisface
\cref{eq:descendant-isotropic-bounce} y ocupa otro estatuto.

\begin{proposition}[Falsabilidad del rebote sincronizado]
\label{prop:descendant-quantization-falsifier}
Fijados prospectivamente \((s_0,\rho_0,w)\), la realizaci\'on sincronizada
existe \'unicamente si el n\'umero de
\cref{eq:descendant-quantized-bounce} es un entero. La distancia
\begin{equation}
 \delta_9:=\operatorname{dist}(q_b,\Z)
 \label{eq:descendant-quantization-defect}
\end{equation}
es un falsador cuantitativo y se anula exactamente en un nivel permitido.
\end{proposition}

\begin{proof}
La condici\'on \(n=9q\) exige \(q\in\Z\). Sustituir
\cref{eq:descendant-nonadic-scale} en
\cref{eq:descendant-isotropic-bounce} y tomar logaritmos produce la f\'ormula.
La caracterizaci\'on de \(\delta_9\) sigue de la definici\'on de distancia a
los enteros.
\end{proof}

\section{Datos de empalme en la hipersuperficie de rebote}

Sea \(\Sigma_b\) una hipersuperficie espacial que separa una regi\'on
contractiva \(\mathcal M_-\) de una regi\'on expansiva
\(\mathcal M_+\). Denotemos por \(h_{ij}\) la m\'etrica inducida, por
\(n^\mu\) la normal orientada y por
\(\mathring K_{ij}^\pm\) las curvaturas extr\'insecas calculadas con la
conexi\'on de Levi--Civita de cada lado. Adoptamos
\begin{equation}
 [X]:=X^+\big|_{\Sigma_b}-X^-\big|_{\Sigma_b},
 \qquad
 \kappa:=\frac{8\pi G}{c^4}.
 \label{eq:descendant-jump-kappa}
\end{equation}

Las condiciones m\'etricas de Israel son
\begin{align}
 [h_{ij}]&=0,
 \label{eq:descendant-Israel-first}\\
 [\mathring K_{ij}-h_{ij}\mathring K]
 &=-\kappa S^{\Sigma}_{ij},
 \label{eq:descendant-Israel-second}
\end{align}
donde \(S^\Sigma_{ij}\) es el tensor energ\'ia--impulso superficial en la
convenci\'on de normal fijada. Un empalme sin capa material satisface
\begin{equation}
 [h_{ij}]=0,
 \qquad
 [\mathring K_{ij}]=0.
 \label{eq:descendant-shell-free-Israel}
\end{equation}
La inversi\'on de \(\Theta\) atraviesa entonces un cero continuo y queda
realizada por la din\'amica, en lugar de introducirse mediante un salto de la
normal.

En la formulaci\'on de primer orden, sea
\(\omega^{ab}=\mathring\omega^{ab}+C^{ab}\) la conexi\'on con contorsi\'on,
y escribamos \(T_\mu:=T^\lambda{}_{\mu\lambda}\) para la traza de la
torsi\'on.
Fijamos la normalizaci\'on de la corriente de esp\'in mediante la ecuaci\'on
de Cartan
\begin{equation}
 \mathscr C(T)^{\lambda}{}_{\mu\nu}
 :=T^{\lambda}{}_{\mu\nu}
 +\delta^\lambda_\mu T_\nu
 -\delta^\lambda_\nu T_\mu
 =\kappa s^{\lambda}{}_{\mu\nu}.
 \label{eq:descendant-Cartan-equation}
\end{equation}
Su parte distribucional sobre \(\Sigma_b\) exige
\begin{equation}
 \boxed{
 \mathscr C(T_\Sigma)
 =\kappa s_\Sigma.}
 \label{eq:descendant-Cartan-junction}
\end{equation}
En ausencia de una capa superficial de esp\'in,
\(s_\Sigma=0\), la conexi\'on admisible carece de una componente torsional
delta y la contorsi\'on tangencial posee un empalme regular. Una acci\'on con
t\'ermino de Holst o acoplamientos fermi\'onicos no m\'inimos reemplaza
\(\mathscr C\) por el operador algebraico correspondiente; las condiciones
se obtienen integrando esa ecuaci\'on a trav\'es de la capa.

\begin{criterion}[Empalme Israel--Cartan regular]
\label{crit:descendant-Israel-Cartan}
Una costura torsional regular a trav\'es de \(\Sigma_b\) requiere:
\begin{enumerate}[label=\roman*)]
 \item continuidad de la coestructura inducida, equivalente a
       \cref{eq:descendant-Israel-first} hasta gauge de Lorentz;
 \item ecuaci\'on de Israel \cref{eq:descendant-Israel-second} con toda capa
       superficial declarada;
 \item ecuaci\'on de Cartan \cref{eq:descendant-Cartan-junction} con toda
       corriente superficial declarada;
 \item continuidad de las cargas de gauge y de la orientaci\'on de ruta;
 \item finitud de los invariantes de curvatura y torsi\'on.
\end{enumerate}
\end{criterion}

La conexi\'on de Ashtekar--Barbero
\begin{equation}
 A_i^a=\Gamma_i^a+\gamma_{\rm BI}K_i^a
 \label{eq:descendant-Barbero-connection}
\end{equation}
puede actuar como lectora de la costura cuando ambos lados comparten
coestructura, \(\gamma_{\rm BI}\) y normalizaci\'on. Su continuidad es una
consecuencia condicionada de los datos de empalme correspondientes; su
verificaci\'on exige transportar expl\'icitamente \(\Gamma\) y \(K\) a la
misma secci\'on de borde.

\begin{proposition}[Involuci\'on conjunta de Barbero y curvatura extr\'inseca]
\label{prop:descendant-Barbero-involution}
Consid\'erese la involuci\'on algebraica
\begin{equation}
 \mathfrak J_b:
 (\Gamma_i^a,K_i^a,\gamma_{\rm BI})
 \longmapsto
 (\Gamma_i^a,-K_i^a,-\gamma_{\rm BI}).
 \label{eq:descendant-Barbero-involution}
\end{equation}
Entonces
\begin{equation}
 \mathfrak J_b(A_i^a)=A_i^a,
 \label{eq:descendant-Barbero-A-invariant}
\end{equation}
y todo espectro areal que dependa de \(|\gamma_{\rm BI}|\) permanece
invariante. La identificaci\'on de \(\mathfrak J_b\) con la inversi\'on
din\'amica del rebote requiere que el entrelazador Israel--Cartan transporte
la orientaci\'on de \(K\) y la hoja de \(\gamma_{\rm BI}\) conforme a
\cref{eq:descendant-Barbero-involution}.
\end{proposition}

\begin{proof}
El producto de los dos factores orientados satisface
\((-\gamma_{\rm BI})(-K_i^a)=\gamma_{\rm BI}K_i^a\), mientras \(\Gamma_i^a\)
queda fija. La dependencia areal respecto de \(|\gamma_{\rm BI}|\) es par.
\end{proof}

\section{Horizonte nulo como hipersuperficie marginal y condición inicial de empalme}

La hipersuperficie del horizonte \(\Sigma_H\) posee car\'acter nulo y requiere
datos de empalme adaptados. Sean \(q_{AB}\) la m\'etrica bidimensional de las
secciones, \(\ell^\mu\) el generador nulo y \(N^\mu\) un transversal con
\(\ell\cdot N=-1\). La curvatura transversal es
\begin{equation}
 \mathcal C_{AB}
 :=q_A{}^\mu q_B{}^\nu\nabla_\mu N_\nu.
 \label{eq:descendant-null-transverse-curvature}
\end{equation}
Un empalme nulo sin capa impulsiva conserva \(q_{AB}\), la clase de
\(\mathcal C_{AB}\) que determina la tensi\'on superficial y el flujo de
Cartan tangencial. La marginalidad
\begin{equation}
 \theta_{(\ell)}\big|_{\Sigma_H}=0
 \label{eq:descendant-null-marginality}
\end{equation}
realiza el primer estado \(0\) de
\cref{eq:descendant-five-stage-sequence}. La formulaci\'on nula mantiene
separada esta condici\'on de la condici\'on volum\'etrica
\(\Theta|_{\Sigma_b}=0\).

\section{Banda de M\"obius y memoria de orientaci\'on}

Sobre las palabras dodecaf\'asicas act\'ua la involuci\'on geneal\'ogica
\begin{equation}
 C_M(\mathbf t):=-\operatorname{rev}(\mathbf t),
 \qquad
 C_M^2=I.
 \label{eq:descendant-CM}
\end{equation}
La l\'inea orientada asociada se obtiene por el cociente
\begin{equation}
 (\mathbf t,\chi)\sim(C_M\mathbf t,-\chi).
 \label{eq:descendant-Mobius-quotient}
\end{equation}
Cuando la base es un lazo de retorno, esta identificaci\'on define una banda
con holonom\'ia \(-I\): una vuelta invierte la orientaci\'on y dos vueltas la
restauran.

Sea \(\mathscr H_\gamma\) la holonom\'ia sobre una vuelta. La jerarqu\'ia
\begin{equation}
 \mathcal R_{\rm vec}
 \subseteq\mathcal R_{\rm ray}
 \subseteq\mathcal R_{\rm obs}
 \label{eq:descendant-return-hierarchy}
\end{equation}
distingue retorno de vector, retorno de rayo y retorno de observable. Para
\(\mathscr H_\gamma=-I\), el rayo retorna en una vuelta, el vector orientado
retorna en dos y la memoria de signo conserva el recorrido.

\begin{proposition}[Estatuto topol\'ogico de la banda]
\label{prop:descendant-Mobius-status}
La relaci\'on \cref{eq:descendant-Mobius-quotient} define una l\'inea real
no orientable sobre el lazo exactamente cuando su primera clase de
Stiefel--Whitney satisface
\begin{equation}
 w_1=1\in H^1(S^1;\Z_2).
 \label{eq:descendant-w1}
\end{equation}
Esta afirmaci\'on concierne al haz geneal\'ogico de orientaci\'on. Una
realizaci\'on como topolog\'ia del espacio-tiempo requiere un morfismo hacia
el haz de marcos o de espinores que conserve causalidad y estructura de
Cartan.
\end{proposition}

\begin{proof}
Una l\'inea real sobre \(S^1\) est\'a determinada por la paridad de su
funci\'on de transici\'on. La transici\'on \(-1\) de
\cref{eq:descendant-Mobius-quotient} representa la clase no trivial; dos
vueltas elevan la transici\'on al cuadrado y producen \(+1\).
\end{proof}

La banda proporciona una representaci\'on exacta de orientaci\'on y memoria.
La continuidad f\'isica a trav\'es del horizonte y del rebote se formula
mediante los empalmes de las secciones anteriores; ambas estructuras se
coordinan mediante un entrelazador y conservan sus dominios.

\section{Observador relativo, transporte paralelo y autoinercia en la región cosmológica descendiente}

El estado din\'amico ampliado incluye observador y memoria:
\begin{equation}
 \Gamma=(x,\mathcal O,m),
 \label{eq:descendant-observer-state}
\end{equation}
y el marco accesible al observador es una restricci\'on de la conexi\'on total,
\begin{equation}
 \mathfrak F_{\mathcal O}(\Gamma)
 :=\operatorname{Res}_{\mathcal O}
 (\Gamma,\nabla_{\TPK}).
 \label{eq:descendant-observer-frame}
\end{equation}
Una trayectoria horizontal del sistema completo,
\begin{equation}
 \nabla^{\rm tot}_{\dot\Gamma}\dot\Gamma=0,
 \label{eq:descendant-total-horizontal}
\end{equation}
puede publicar aceleraci\'on en un subsistema \(S\):
\begin{equation}
 a_S=\nabla^S_{\dot\gamma_S}\dot\gamma_S
 =\mathcal K_{\pi_S}(\dot\Gamma,\dot\Gamma),
 \label{eq:descendant-projected-acceleration}
\end{equation}
donde \(\mathcal K_{\pi_S}\) es el segundo tensor fundamental de la
proyecci\'on. La autoinercia designa la horizontalidad de la totalidad y la
respuesta inducida en sus proyecciones.

El acoplamiento entre sistema, aparato y registro se describe mediante una
isometr\'ia
\begin{equation}
 V_{\rm obs}:\mathcal H_{\rm in}\longrightarrow
 \mathcal H_{\rm sys}\otimes\mathcal H_{\rm mem},
 \qquad
 V_{\rm obs}^*V_{\rm obs}=I.
 \label{eq:descendant-observer-isometry}
\end{equation}
Sea \(\{P_y\}\) la PVM del registro en \(\mathcal H_{\rm mem}\), con
\begin{equation}
 P_yP_z=\delta_{yz}P_y,
 \qquad
 \sum_yP_y=I_{\rm mem}.
 \label{eq:descendant-observer-PVM}
\end{equation}
La isometr\'ia y esta PVM inducen el instrumento completamente positivo
\begin{equation}
 \mathcal I_y(\rho)
 :=\Tr_{\rm mem}\!\left[
 (I_{\rm sys}\otimes P_y)V_{\rm obs}\rho V_{\rm obs}^*
 (I_{\rm sys}\otimes P_y)
 \right],
 \qquad
 \sum_y\Tr\mathcal I_y(\rho)=\Tr\rho.
 \label{eq:descendant-observer-instrument}
\end{equation}
Los estados relativos se obtienen entonces mediante
\begin{equation}
 p(y)=\Tr\mathcal I_y(\rho),
 \qquad
 \rho_y=\frac{\mathcal I_y(\rho)}{p(y)}.
 \label{eq:descendant-relative-state}
\end{equation}
La secci\'on \(\rho_y\) es el estado relativo al registro \(y\); el estado
global conserva todos los canales en la imagen de la isometr\'ia. Esta
estructura formaliza el parentesco con Everett en el nivel preciso de
estados relativos y deja la ontolog\'ia de las ramas como elecci\'on f\'isica
posterior.

\section{Canal global de la regi\'on descendiente}

La versi\'on operatoria de una apertura descendiente utiliza una isometr\'ia
\begin{equation}
 V_{\rm pc}:\mathcal H_{\rm in}\longrightarrow
 \mathcal H_{\rm parent}\otimes\mathcal H_{\rm child},
 \qquad
 V_{\rm pc}^*V_{\rm pc}=I.
 \label{eq:descendant-parent-child-isometry}
\end{equation}
El canal conjunto y sus canales complementarios son
\begin{align}
 \mathcal N_{\rm pc}(\rho)
 &=V_{\rm pc}\rho V_{\rm pc}^*,
 \label{eq:descendant-global-channel}\\
 \mathcal N_{\rm parent}(\rho)
 &=\Tr_{\rm child}\mathcal N_{\rm pc}(\rho),
 \qquad
 \mathcal N_{\rm child}(\rho)
 =\Tr_{\rm parent}\mathcal N_{\rm pc}(\rho).
 \label{eq:descendant-complementary-channels}
\end{align}

\begin{theorem}[Conservaci\'on global y distribuci\'on relativa]
\label{thm:descendant-global-conservation}
La aplicaci\'on \(\mathcal N_{\rm pc}\) conserva traza y pureza de un estado
puro sobre su imagen. Los canales parentales y descendientes pueden producir
estados mixtos, y para una salida bipartita pura satisfacen
\begin{equation}
 S(\rho_{\rm parent})=S(\rho_{\rm child}).
 \label{eq:descendant-equal-entanglement}
\end{equation}
La informaci\'on global queda contenida en las salidas y sus correlaciones.
\end{theorem}

\begin{proof}
La identidad \(V_{\rm pc}^*V_{\rm pc}=I\) implica
\[
 \Tr(V_{\rm pc}\rho V_{\rm pc}^*)
 =\Tr(\rho V_{\rm pc}^*V_{\rm pc})=\Tr\rho.
\]
Una isometr\'ia lleva vectores unitarios a vectores unitarios. La igualdad de
entrop\'ias reducidas sigue de la descomposici\'on de Schmidt del estado
bipartito puro.
\end{proof}

La conservaci\'on de una carga \(Q\) o de la energ\'ia exige adem\'as los
entrelazamientos
\begin{equation}
 Q_{\rm out}V_{\rm pc}=V_{\rm pc}Q_{\rm in},
 \qquad
 H_{\rm out}V_{\rm pc}=V_{\rm pc}H_{\rm in}.
 \label{eq:descendant-charge-energy-intertwining}
\end{equation}
La isometr\'ia conserva probabilidad; las ecuaciones
\cref{eq:descendant-charge-energy-intertwining} conservan las magnitudes
din\'amicas declaradas. Una salida admisible distribuye la informaci\'on en
correlaciones y respeta la imposibilidad de clonar un estado arbitrario en
dos copias id\'enticas.

La procedencia de los sectores \(90/120\) se conserva mediante el
Hamiltoniano hologr\'afico modular
\begin{equation}
 K_{\rm HHM}=c_0I+\Delta_{\rm mod}
 (90N_{90}+120N_{120}).
 \label{eq:descendant-HHM}
\end{equation}
La ocupaci\'on total y la memoria ponderada,
\begin{equation}
 N=N_{90}+N_{120},
 \qquad
 D=90N_{90}+120N_{120},
 \label{eq:descendant-HHM-data}
\end{equation}
El transporte de la ocupaci\'on total se impone como identidad operatoria
independiente:
\begin{equation}
 N^{\rm out}V_{\rm pc}=V_{\rm pc}N^{\rm in}.
 \label{eq:descendant-number-intertwining}
\end{equation}
Conservados \(N\) y la memoria ponderada \(D\), los dos sectores se
reconstruyen un\'ivocamente:
\begin{equation}
 N_{90}=4N-\frac D{30},
 \qquad
 N_{120}=\frac D{30}-3N.
 \label{eq:descendant-HHM-recovery}
\end{equation}
La condici\'on de transporte
\begin{equation}
 K_{\rm HHM}^{\rm out}V_{\rm pc}
 =V_{\rm pc}K_{\rm HHM}^{\rm in}
 \label{eq:descendant-HHM-intertwining}
\end{equation}
Las identidades
\cref{eq:descendant-number-intertwining,eq:descendant-HHM-intertwining}
preservan conjuntamente ocupaci\'on y procedencia sectorial a trav\'es de la
costura. Estas igualdades se a\~naden a las obligaciones din\'amicas de una
realizaci\'on completa.

\section{Continuaci\'on descendiente y cobordismo ramificado}

La denominaci\'on \emph{universo descendiente} admite dos niveles
geom\'etricos.

\begin{definition}[Continuaci\'on descendiente m\'inima]
\label{def:descendant-minimal}
Una continuaci\'on descendiente es una regi\'on
\(\mathcal M_{\rm child}\) situada al futuro del rebote, con
\begin{equation}
 \Theta>0,
 \qquad
 \text{curvatura y torsi\'on finitas},
 \qquad
 \text{una funci\'on temporal global en su dominio causal}.
 \label{eq:descendant-minimal-conditions}
\end{equation}
La regi\'on puede pertenecer a una prolongaci\'on conectada del interior y
conserva la topolog\'ia de las hojas espaciales.
\end{definition}

\begin{definition}[Cobordismo parental--descendiente]
\label{def:descendant-cobordism}
La versi\'on ramificada es un cobordismo efectivo
\begin{equation}
 \mathcal W:\Sigma_{\rm in}\longrightarrow
 \Sigma_{\rm parent}\sqcup\Sigma_{\rm child},
 \label{eq:descendant-cobordism}
\end{equation}
equipado con estructura de Cartan, datos de empalme
\cref{crit:descendant-Israel-Cartan} y canal
\cref{eq:descendant-parent-child-isometry}.
\end{definition}

La continuaci\'on m\'inima requiere un rebote y una extensi\'on causal. El
cobordismo ramificado a\~nade cambio topol\'ogico y separaci\'on de salidas.
En una geometr\'ia lorentziana cl\'asica, no degenerada y globalmente
hiperb\'olica, el cambio topol\'ogico est\'a sujeto a obstrucciones causales.
La versi\'on ramificada debe, por ello, realizarse mediante una regi\'on
cu\'antica o torsional donde se declare qu\'e hip\'otesis cl\'asica se modifica,
o mediante una interpretaci\'on de dos dominios causales dentro de una misma
topolog\'ia global.

\begin{criterion}[Promoci\'on a universo descendiente f\'isico]
\label{crit:descendant-promotion}
La imagen completa adquiere estatuto de realizaci\'on f\'isica cuando se
construyen simult\'aneamente:
\begin{enumerate}[label=\roman*)]
 \item una soluci\'on de las ecuaciones Einstein--Cartan con
       \cref{eq:descendant-bounce-conditions};
 \item datos de Israel--Cartan regulares en \(\Sigma_H\) y \(\Sigma_b\);
 \item una regi\'on expansiva futura con funci\'on temporal y superficies de
       Cauchy en cada rama accesible;
 \item finitud de invariantes y prolongaci\'on de geod\'esicas a trav\'es del
       n\'ucleo;
 \item ausencia de curvas causales cerradas y de capas no registradas;
 \item un canal isom\'etrico que conserve las cargas declaradas y el
       Hamiltoniano de procedencia \cref{eq:descendant-HHM-intertwining};
 \item naturalidad non\'adica de la corriente torsional y de los datos de
       frontera;
 \item un entrelazador entre la banda geneal\'ogica y el haz f\'isico de
       marcos o espinores.
\end{enumerate}
\end{criterion}

Estas condiciones transforman la narraci\'on
\cref{eq:descendant-five-stage-sequence} en un problema de existencia
geom\'etrico, operatorio y causal con controles separados.

\section{Consecuencias locales del empalme torsional: conservación, regularidad y prolongación causal}

\begin{proposition}[Contenido de una realizaci\'on completa]
\label{prop:descendant-conditioned-consequences}
Sup\'ongase satisfecho el
\cref{crit:descendant-promotion}. Entonces:
\begin{enumerate}[label=\alph*)]
 \item el signo causal recorre la secuencia
       \cref{eq:descendant-five-stage-sequence};
 \item la memoria torsional conserva la covariancia \(V^{-2}\) a ambos
       lados del rebote;
 \item el estado global evoluciona mediante un canal que conserva traza;
 \item cada observador accede a un estado relativo de su rama;
 \item las cargas incluidas en
       \cref{eq:descendant-charge-energy-intertwining} se conservan;
 \item la orientaci\'on de M\"obius se transporta como dato de holonom\'ia.
\end{enumerate}
\end{proposition}

\begin{proof}
Cada afirmaci\'on es la aplicaci\'on de una de las condiciones del criterio:
el lector causal prueba a); el transductor torsional prueba b); la isometr\'ia
prueba c) y d); los entrelazamientos prueban e); y el morfismo de haces prueba
f).
\end{proof}

La proposici\'on organiza consecuencias de una estructura construida. La
existencia de la continuaci\'on m\'inima se resuelve mediante el
\cref{thm:ivv-descendant-existence}; la ramificaci\'on cl\'asica con cambio
topol\'ogico queda clasificada por la obstrucci\'on del
\cref{thm:ivv-cobordism-obstruction}.

\section{Dinámica torsional del rebote: identidades exactas, interfaces cosmológicas y condiciones de cierre}

\begin{center}
\small
\begin{tabularx}{\textwidth}{>{\bfseries\raggedright\arraybackslash}p{31mm}YY}
\toprule
Objeto & Estatuto & Obligaci\'on o control\\
\midrule
Ley \(M_n=V_n^{-2}\) & Resultado interno exacto & Certificado de Perron y refinamiento\\
Lector \(-\operatorname{sgn}\vartheta\) & Formalizaci\'on exacta & Observable de expansi\'on tipado\\
Rebote homog\'eneo & Consecuencia condicionada & \(s_0>0\), \(w<1\), cero transversal\\
Rebote Kantowski--Sachs & Interfaz f\'isica & Dominancia de torsi\'on, cizalla y curvatura controladas\\
Empalme Israel--Cartan & Interfaz f\'isica & Datos de capa y corriente de esp\'in declarados\\
Banda de M\"obius & Resultado topol\'ogico exacto & Entrelazador hacia el haz f\'isico\\
Estado relativo & Resultado operatorio exacto & Instrumento y fibra observada\\
Canal progenitor--descendiente & Formalizaci\'on condicionada & Isometr\'ia y conservaci\'on de cargas\\
Continuaci\'on descendiente & Teorema de existencia & Soluci\'on regular y causalmente completa\\
Cobordismo ramificado & Obstrucci\'on clásica demostrada & Incompatibilidad con global hiperbólica y cambio topol\'ogico\\
\bottomrule
\end{tabularx}
\end{center}

\section{Obstrucciones a la continuación cosmológica: pérdida de orientación, regularidad del empalme o memoria de retorno}

Los controles negativos distinguen cada capa:
\begin{enumerate}
 \item una soluci\'on con \(\Theta=0\) y \(\dot\Theta<0\) representa un
       m\'aximo de volumen y queda fuera del rebote;
 \item el caso \(8\pi Gs_0\leq\Sigma_0^2\) elimina la garant\'ia de
       dominancia torsional anis\'otropa;
 \item un \(q_b\notin\Z\) conserva el rebote continuo y refuta la versi\'on
       sincronizada con una vuelta completa;
 \item un salto de \(h_{ij}\), de la cantidad de Israel o de la corriente de
       Cartan requiere una capa superficial expl\'icita;
 \item una banda con \(w_1=0\) posee orientaci\'on global y refuta la
       realizaci\'on de M\"obius;
 \item un mapa \(V_{\rm pc}\) con \(V_{\rm pc}^*V_{\rm pc}\ne I\) pierde
       probabilidad global;
 \item dos copias id\'enticas de un estado arbitrario violan el control de no
       clonaci\'on y refutan la lectura de distribuci\'on correlacionada;
 \item una regi\'on descendiente con curvas causales cerradas, invariantes
       divergentes o geod\'esicas inextendibles incumple la promoci\'on causal;
 \item un cambio topol\'ogico que conserve simult\'aneamente todas las
       hip\'otesis cl\'asicas obstruidas carece del dominio geom\'etrico
       requerido.
\end{enumerate}

\begin{proofstatusbox}
La gram\'atica \(-1,0,+1\), el lector de signo, la involuci\'on
\(C_M^2=I\), la clasificaci\'on de la banda, la ley \(V^{-2}\), la jerarqu\'ia
de retornos y la conservaci\'on de un canal isom\'etrico son resultados
exactos en sus dominios. El rebote Einstein--Cartan es una consecuencia de
las ecuaciones y desigualdades declaradas. La extensi\'on a
Kantowski--Sachs, los empalmes, la continuidad de la conexi\'on de Barbero y
la apertura causal son interfaces f\'isicas. La continuaci\'on
parental--descendiente m\'inima queda realizada por el
\cref{thm:ivv-descendant-existence}; el cambio topol\'ogico ramificado bajo
las hip\'otesis cl\'asicas queda excluido por el
\cref{thm:ivv-cobordism-obstruction}.
\end{proofstatusbox}

\begin{falsifierbox}
La costura torsional queda refutada si la corriente construida desde HMT
produce una potencia distinta de \(V^{-2}\), una amplitud de signo
incompatible o un fallo de refinamiento. El rebote queda refutado por la
ausencia de una ra\'iz transversal regular. El empalme queda refutado por un
residuo de Israel o Cartan sin fuente superficial. La conservaci\'on global
queda refutada por p\'erdida de traza, de norma o de una carga entrelazada.
La promoci\'on a universo descendiente queda refutada por singularidad,
violaci\'on causal o imposibilidad de construir el haz de orientaci\'on y el
canal global exigidos. Cada refutaci\'on afecta a su interfaz y conserva los
resultados internos anteriores.
\end{falsifierbox}
